\section{Introduction}
An agent at a pickup kiosk can read an order status, a nearby sign and an
attacker's overlay through the same camera. Correctly decoding a machine-readable
marker does not establish who issued its claim, which interaction it belongs to,
or whether its owner authorized the agent's proposed action. Even a valid signed
response can be copied to another display or relayed from another location.
Treating the entire camera view as authoritative because one fragment verifies
therefore strengthens the evidence without justification.

OpenAST (Open Agent Sight Trust) implements a narrow alternative. A participating
display signs the exact requested statement and a challenge bound to the agent's
complete proposed action. The signed evidence travels through an actual QR raster
and is decoded independently of surrounding text. A separate owner grant and
trusted gateway determine whether that action can create a durable local receipt.
A successful exchange supports issuer-key control, message integrity, a fresh
signed response and exact request binding. It does not establish scene truth,
physical location, proximity or consent.

This separation is useful when an agent must use a human-facing display rather
than a compatible authenticated API. If an appropriate API already supplies the
claim and permission, it may be the better integration. OpenAST is not a reason
to replace a sufficient API with a QR channel. Its scope is an inspectable
reference for the interface-to-action boundary.

\paragraph{What this report establishes.}
We provide a runnable Python visual adapter, shared permission gate, real QR
encoding/decoding, finite attack cases, generated-image reliability trials and
explicit negative results. We compare with both a signature-only ablation and a
control implementing the same serial policy. Agreement with the stronger control
limits the novelty claim: no additional authorization property is demonstrated.
The open research question is how scoped evidence can later compose across voice,
sight and touch while avoiding unnecessary acquisition. That larger mechanism
is not implemented or proved here.

